% answers.tex —— 1955 年·全国卷 · 参考答案与详细解析片段
% 由 tools/generate_projects.py 依来源总账自 1955/1955_数学真题_试题.md + 1955/1955_数学真题_解析正文.md + 1955/1955_数学真题_答案要点.md 生成，勿手改（改源 MD 后重跑）
\gksubpart{一、小题}

\begin{answer}
\gksolhead{1}{}

\textbf{题目}：

（甲）以二次方程 $x^2-3x-1=0$ 的两根的平方为两根作一二次方程．

\Needspace{4\baselineskip}
\textbf{解析}

设 $x^2-3x-1=0$ 的两根为 $\alpha,\beta$，由韦达定理

\nopagebreak[4]
\[
\alpha+\beta=3,\qquad \alpha\beta=-1.
\]

两根的平方之和与积：

\nopagebreak[4]
\[
\alpha^2+\beta^2=(\alpha+\beta)^2-2\alpha\beta=9-2(-1)=11,
\]

\nopagebreak[4]
\[
\alpha^2\beta^2=(\alpha\beta)^2=1.
\]

故以 $\alpha^2,\beta^2$ 为根的二次方程为 $y^2-11y+1=0$。

\gkanswer{答案：$y^2-11y+1=0$}

\end{answer}
\begin{answer}
\gksolhead{2}{}

\textbf{题目}：

（乙）等腰三角形一腰的长是底边的 $4$ 倍，求这三角形各角的余弦．

\Needspace{4\baselineskip}
\textbf{解析}

设 $\triangle ABC$ 中 $AB=AC=4BC$。令 $BC=x$，则 $AB=AC=4x$。由余弦定理

\nopagebreak[4]
\[
\cos A=\frac{AB^2+AC^2-BC^2}{2\,AB\cdot AC}
=\frac{16x^2+16x^2-x^2}{2\cdot4x\cdot4x}
=\frac{31x^2}{32x^2}=\frac{31}{32}.
\]

取 $BC$ 边上的高 $AD$，则 $D$ 是 $BC$ 中点，$CD=\dfrac x2$。在直角三角形 $ADC$ 中

\nopagebreak[4]
\[
\cos C=\frac{CD}{AC}=\frac{x/2}{4x}=\frac18,
\]

由等腰知 $\cos B=\cos C=\dfrac18$。

\gkanswer{答案：$\cos A=\dfrac{31}{32}$，$\cos B=\cos C=\dfrac18$}

\end{answer}
\begin{answer}
\gksolhead{3}{}

\textbf{题目}：

（丙）已知正四棱锥底边的长为 $a$，侧棱与底面的交角为 $45^\circ$，求这棱锥的高．

\Needspace{4\baselineskip}
\textbf{解析}

设正四棱锥 $S-ABCD$，高为 $SH$，底边长 $a$，侧棱与底面的交角为 $\angle SAH=45^\circ$。底面正方形对角线的一半

\nopagebreak[4]
\[
AH=\frac{\sqrt2}{2}a.
\]

在 $\triangle SHA$ 中，$\angle SHA=90^\circ$，$\angle SAH=45^\circ$，故 $\triangle SHA$ 为等腰直角三角形，$SH=AH$。因此

\nopagebreak[4]
\[
SH=\frac{\sqrt2}{2}a.
\]

\gkanswer{答案：$\dfrac{\sqrt2}{2}a$}

\end{answer}
\begin{answer}
\gksolhead{4}{}

\textbf{题目}：

（丁）写出：二面角的平面角的定义．

\Needspace{4\baselineskip}
\textbf{解析}

二面角的平面角的定义：自二面角的棱上一点，在它的两个面内分别作棱的垂线，这两条垂线所夹的角，叫做这个二面角的平面角。

\gkanswer{答案：自二面角的棱上一点，在其两个面内分别作棱的垂线，这两条垂线所夹的角叫做该二面角的平面角。}

\par\medskip

\gksubpart{二、大题}

\end{answer}
\begin{answer}
\gksolhead{5}{}

\textbf{题目}：

（二）求 $b$，$c$，$d$ 的值，使多项式 $x^3+bx^2+cx+d$ 适合下列三条件：（1）被 $x-1$ 整除；（2）被 $x-3$ 除时余 $2$；（3）被 $x+2$ 除与被 $x-2$ 除时余数相等．

\Needspace{4\baselineskip}
\textbf{解析}

记 $f(x)=x^3+bx^2+cx+d$。由三条件：

（1）$f(x)$ 被 $x-1$ 整除，即 $f(1)=0$：

\nopagebreak[4]
\[
1+b+c+d=0.
\]

（2）$f(x)$ 被 $x-3$ 除余 $2$，即 $f(3)=2$：

\nopagebreak[4]
\[
27+9b+3c+d=2.
\]

（3）被 $x+2$ 除与被 $x-2$ 除余数相等，即 $f(-2)=f(2)$：

\nopagebreak[4]
\[
-8+4b-2c+d=8+4b+2c+d
\implies -4c=16\implies c=-4.
\]

代入前两式：$b+d=3$，$9b+d=-13$。两式相减 $8b=-16$，故 $b=-2$，进而 $d=5$。

\gkanswer{答案：$b=-2$，$c=-4$，$d=5$}

\end{answer}
\begin{answer}
\gksolhead{6}{}

\textbf{题目}：

（三）由直角三角形勾上一点 $D$ 作弦 $AB$ 的垂线交弦于 $E$、股的延长线于 $F$、外接圆周于 $Q$，求证：$EQ$ 为 $EA$ 与 $EB$ 的比例中项又为 $ED$ 与 $EF$ 的比例中项．

\Needspace{4\baselineskip}
\textbf{解析}

连接 $QA,QB$。因 $Q$ 在直角三角形的外接圆上，$AB$ 为直径，故 $\angle AQB=90^\circ$。又 $EQ\perp AB$，故 $EQ$ 是直角三角形 $AQB$ 斜边 $AB$ 上的高，于是

\nopagebreak[4]
\[
EQ^2=EA\cdot EB.
\]

又 $\angle 1=\angle 2$（它们都与 $\angle ABF$ 互余），且 $\angle AED=\angle FEB=90^\circ$，故 $\triangle AED\sim\triangle FEB$，得

\nopagebreak[4]
\[
\frac{EA}{EF}=\frac{ED}{EB}
\implies EA\cdot EB=ED\cdot EF.
\]

综合得 $EQ^2=EA\cdot EB=ED\cdot EF$，即 $EQ$ 既是 $EA,EB$ 的比例中项，又是 $ED,EF$ 的比例中项。

\gkanswer{答案：见解析版}

\end{answer}
\begin{answer}
\gksolhead{7}{}

\textbf{题目}：

（四）解方程 $\cos2x=\cos x+\sin x$，求 $x$ 的通值．

\Needspace{4\baselineskip}
\textbf{解析}

\nopagebreak[4]
\[
\cos2x=\cos^2x-\sin^2x=(\cos x+\sin x)(\cos x-\sin x).
\]

原方程化为

\nopagebreak[4]
\[
(\cos x+\sin x)(\cos x-\sin x)=\cos x+\sin x,
\]

\nopagebreak[4]
\[
(\cos x+\sin x)(\cos x-\sin x-1)=0.
\]

由 $\cos x+\sin x=0$ 得 $\tan x=-1$，$x=n\pi-\dfrac\pi4$。

由 $\cos x-\sin x-1=0$，即 $\cos x-\sin x=1$，两边同乘 $\dfrac{\sqrt2}{2}$：

\nopagebreak[4]
\[
\cos\left(x+\frac\pi4\right)=\frac{\sqrt2}{2}=\cos\frac\pi4,
\]

故 $x+\dfrac\pi4=2n\pi\pm\dfrac\pi4$，即 $x=2n\pi$ 或 $x=2n\pi-\dfrac\pi2$。

\gkanswer{答案：$x=n\pi-\dfrac\pi4$，$x=2n\pi$，$x=2n\pi-\dfrac\pi2$（$n$ 为整数）}

\end{answer}
\begin{answer}
\gksolhead{8}{}

\textbf{题目}：

（五）一三角形三边的长成等差级数，其周长为 $12$ 尺，面积为 $6$ 平方尺，求证这三角形为一直角三角形．

\Needspace{4\baselineskip}
\textbf{解析}

设三角形三边为 $x-y,x,x+y$（$x>y>0$），成等差数列。周长为 $12$：

\nopagebreak[4]
\[
(x-y)+x+(x+y)=12\implies 3x=12\implies x=4.
\]

半周长 $p=6$。由海伦公式，面积

\nopagebreak[4]
\[
S=\sqrt{6(6-x-y)(6-x)(6-x+y)}.
\]

代入 $x=4$、$S=6$ 并平方：

\nopagebreak[4]
\[
36=6(2-y)(2)(2+y)
\implies 12(2-y)(2+y)=36
\implies 4-y^2=3
\implies y^2=1.
\]

取 $y=1$，三边为 $3,4,5$。因 $3^2+4^2=5^2$，由勾股定理逆定理知该三角形为直角三角形。

\gkanswer{答案：见解析版}

\end{answer}